\subsection{RQ1: Where LLM Assistance Contributes}\label{sec:rq1}

RQ1 is best answered as a gradient rather than a single verdict: autonomy is highest for the bounded, well-specified per-module tasks and falls away for the open-ended, system-level work of integration.

Architecture mapping is the one stage that proved not to need an LLM. Its operations (graph tracing of the quantised model, batch-normalisation folding, and weight emission) are deterministic, and in the deployed pipeline they run as a deterministic ONNX frontend and a \texttt{read\_weights} tool rather than a model; the designs built on the resulting layer description pass the byte-exact gate, confirming the extraction is correct without an LLM.

RTL generation, performed by the Foundry, is where LLM assistance contributes most, and its main advantage is generality. One Foundry, with no per-architecture changes and no retraining, generated the per-layer RTL for three networks spanning two architectural families (two ResNet scales and a depthwise-separable MobileNetV2), all of which pass the assembled network byte-exact gate, directly from post-training quantisation and without the per-network quantisation-aware retraining the deterministic flows require. Generation was reliable as well as general: across the ResNet-50 backbone the pipeline averaged 1.09 attempts per state entry.

Closed-loop repair, performed by the Surgeon, is valuable but situational. About 91\% of state entries completed after a single Foundry attempt with no Surgeon intervention; MobileNetV2 drove more repair (1.30 attempts per entry), reflecting harder layer types such as its depthwise and ReLU6 operators. Repair was effective where invoked: ResNet-50 produced zero fail-aborts, so every module that escalated to the Surgeon was resolved.

The pipeline also supports objective-directed refinement of verified RTL: an improve operation supplies a verified module and its Vivado synthesis report to an LLM under one of seven predefined PPA objectives, reverifies each candidate with Verilator, and accepts it only when a deterministic objective-specific threshold is met. Representative per-module runs raised DSP inference and cut LUT and flip-flop use by large margins with correctness preserved (Table~\ref{tab:improve}); these are isolated demonstrations, not claimed to improve the final routed networks.

Orchestration and per-module verification are deterministic by design: the TypeScript orchestrator and the static C++ checker were never delegated to a model, keeping control flow, retry limits, and synthesis dispatch reproducible and the checker outside the generative loop to avoid the two-bug problem (Section~\ref{sec:verification}).

Whole-network integration is the human-directed end of the gradient: System~2 (Section~\ref{sec:system2}) performed all coding, debugging, and tool invocation, while research goals, milestones, and design forks stayed with the author.

Across all System~1 agents, the recorded LLM cost was roughly one to two US dollars per passing module (Table~\ref{tab:rq1cost}), excluding the separate System~2 integration agent.

In short, RQ1's answer is a gradient: autonomy tracks how bounded a task is, near-autonomous on per-module generation and human-directed on integration. Because we did not evaluate deterministic replacements or role-removal ablations for the generative stages, these results characterise the observed effectiveness of those stages rather than the causal marginal benefit of the LLM at each.
